\chapter{TARTAN: Trajectory-Aware Recursive Tiling with Annotated Noise}
\label{app:tartan}

\noindent \TARTAN is a computational framework for multiscale simulation and
inference that distinguishes local plausibility from global consistency using
sheaf-theoretic ideas, treating hallucination as a gluing failure rather than
a noise artifact. This appendix develops the architecture, annotated noise
schema, and four projection modes.

\section{Core Architecture}

\begin{definition}{TARTAN Tiling}
A \emph{TARTAN tiling} of a computational domain $\Omega$ is a hierarchical
partition $\mathcal{T} = \{T_i^{(k)}\}_{i,k}$ where $k$ indexes resolution
levels and $T_i^{(k)}$ are tiles at level $k$. The tiling satisfies:
\begin{itemize}
  \item \textbf{Recursion:} Each tile $T_i^{(k)}$ is covered by a collection
    of tiles $\{T_j^{(k+1)}\}$ at the next resolution level.
  \item \textbf{Trajectory awareness:} Each tile carries a trajectory annotation
    $\tau_i^{(k)}$: the record of prior state sequences that have passed through
    the tile's domain.
  \item \textbf{Noise annotation:} Each tile carries a noise annotation
    $\eta_i^{(k)}$: the record of unresolved structure within the tile that
    could not be assigned to any trajectory.
\end{itemize}
\end{definition}

The fundamental departure from standard multi-resolution analysis is the noise
annotation. Rather than discarding unresolved structure as error to be minimized,
\TARTAN archives it as annotated residue available for future reinterpretation.
This operationalizes the hallucination-as-normal principle of
Chapter~\ref{ch:vocabulary}: the unresolved structure is the non-trivial $H^1$
cocycle of the local sheaf.

\section{The Annotated Noise Schema}

Each noise annotation $\eta_i^{(k)}$ carries the following metadata:

\begin{itemize}
  \item \textbf{Origin:} Which trajectory segment produced this unresolved
    structure.
  \item \textbf{Constraint violation type:} Which admissibility condition the
    structure failed (thermodynamic, logical, geometric, historical).
  \item \textbf{Reinterpretation candidates:} The set of alternative
    trajectory segments that could absorb this structure under relaxed
    gluing conditions.
  \item \textbf{Persistence weight:} How long the noise has remained unresolved;
    high-persistence noise is prioritized for repair.
\end{itemize}

The annotated noise schema converts gluing failures from silent errors into
structured records. A \TARTAN system operating over a long time horizon
accumulates an annotated noise ledger that functions as a diagnostic tool:
recurring noise patterns indicate persistent admissibility failures that
require structural repair rather than local patching.

\section{Local Plausibility versus Global Consistency}

\claimH The central insight of \TARTAN is that local plausibility and global
consistency are distinct properties that standard multi-resolution methods
conflate. A locally plausible configuration is one that satisfies the
admissibility conditions within a single tile. A globally consistent configuration
is one whose local solutions glue into a coherent global section.

\begin{definition}{TARTAN Consistency Check}
A tiling $\mathcal{T}$ is \emph{globally consistent at level $k$} if for every
pair of adjacent tiles $(T_i^{(k)}, T_j^{(k)})$ sharing a boundary $\partial_{ij}$:
\[
  \tau_i^{(k)}\big|_{\partial_{ij}} = \tau_j^{(k)}\big|_{\partial_{ij}},
\]
that is, the trajectory annotations agree on the shared boundary. A failure of
this condition is a \emph{seam inconsistency}, recorded in the noise annotation
of both tiles.
\end{definition}

\section{Four Projection Modes}

\TARTAN operates in four projection modes corresponding to the four derived
stacks of the \RSVP field system (Appendix~\ref{app:rsvp}):

\textbf{Mode 1: Thermodynamic projection.} Each tile tracks the local entropy
production rate. Seam inconsistencies generate entropy; repairs reduce it.
The tiling evolves by preferentially extending tiles that reduce overall entropy
production.

\textbf{Mode 2: Kinetic projection.} Each tile tracks the local velocity field.
Seam inconsistencies manifest as velocity discontinuities at tile boundaries.
The tiling evolves by smoothing velocity across boundaries while preserving
trajectory annotation.

\textbf{Mode 3: Informational projection.} Each tile tracks the local residue
field $\mathcal{R}$. Seam inconsistencies manifest as conflicting provenance
records at tile boundaries. The tiling evolves by resolving provenance conflicts
through the sheaf repair operations of Chapter~\ref{ch:sheaf}.

\textbf{Mode 4: Semantic projection.} Each tile tracks the local semantic
density $\Phi$ and its annotation. Seam inconsistencies are exactly the
semantic tears of the Yarncrawler framework---mismatches between adjacent
local interpretations. The tiling evolves by the stigmergic repair mechanism
of Theorem~\ref{thm:stigmergic}.

\section{Relationship to Hallucination Formalism}

\claimH Under \TARTAN, hallucination is precisely a seam inconsistency that
has been filled by a locally plausible but globally inconsistent trajectory
annotation. The annotation satisfies the admissibility conditions within its
tile but fails the global consistency check at tile boundaries. The annotated
noise schema records this failure with its reinterpretation candidates, enabling
future repair when additional evidence arrives.

This makes the corrigibility condition of Definition~\ref{def:hallucination}
computationally explicit: a corrigible hallucination is one whose seam
inconsistency is recorded in the noise annotation with reinterpretation
candidates; a drifting hallucination is one whose seam inconsistency has been
silently propagated to adjacent tiles without annotation, corrupting their
trajectory records.
